%% sections/03-components.tex

\section{The Campaign Spine Components}
\label{sec:components}

The campaign spine is a single operational document, maintained as a markdown file
throughout the campaign, with five required sections. The document is written before
Session 1 and updated after every session with campaign-permanent facts --- outcomes
that cannot be undone by future sessions --- and any amendments to pending adventure
designs. No amendment to the central question or end-conditions was made after
Session 2 in any of the 7 campaigns, suggesting these components reach early stability
once the campaign's narrative direction is established.

\subsection{Component 1: Central Question}

The central question is a single sentence that the campaign's events will answer by
the finale. Design requirements: (a)~it must be answerable \textit{in principle} by
party behavior; (b)~it must be emotionally resonant for the party archetype in play;
and (c)~it must not be answerable by the module structure alone --- the answer
requires party choices. Examples from the corpus:

\begin{itemize}
  \item \textbf{C1 Moon-Silver Cycle:} \textit{Can grief be returned to the place
    where it began and completed?}
  \item \textbf{C2 Conclave Compact:} \textit{When a covenant breaks, who bears the
    cost of mending it?}
  \item \textbf{C3 Halted Spire:} \textit{Who gets to decide who is worth saving,
    and who gets to decide what counts as faith?}
  \item \textbf{C4 Thorngate Watch:} \textit{How much of yourself do you spend on
    people who don't trust you yet?}
  \item \textbf{C7 Strangers:} \textit{When people with nothing in common have to
    act as one, what do they actually share?}
\end{itemize}

The central question is not a plot question (``who killed the duke?'') but an
existential-choice question (``what are you willing to pay?''). It cannot be
answered by information or by module sequence; it is answered by accumulated
party actions that demonstrate a position on the question. The campaign spine
treats the central question as the organizing principle for all five components:
hints advance it, routes answer it at different depths, spotlights test individual
characters against it, and the deliverable is its final articulation.

\subsection{Component 2: Four-Hint Delivery Schedule}

Four hints are planted in the spine document before the campaign begins and
delivered across sessions. Each hint advances the central question from a
different angle, building toward the finale's simultaneous convergence event.
The designed delivery schedule across a 7-session campaign is:

\begin{itemize}
  \item \textbf{Hint 1 (Sessions 1--2):} Introduces a connection or artifact
    whose significance the party does not yet understand. Designed for
    auto-delivery at passive Perception 15+. Requires no party action.

  \item \textbf{Hint 2 (Sessions 2--4):} Opens a structural or historical
    fact that changes how the party understands the campaign's central problem.
    Often delivered through investigation, NPC dialogue, or a secondary
    artifact discovery.

  \item \textbf{Hint 3 (Sessions 5--6):} Delivered when accumulated party
    engagement unlocks it --- frequently tied to an individual PC's spotlight
    arc. Requires the party to have engaged with the campaign's material
    at a minimum threshold.

  \item \textbf{Hint 4 (Sessions 6--7):} The final piece enabling the
    optimal route. Designed for auto-delivery at the finale scene at a
    passive threshold (Investigation 15 in most campaigns), ensuring
    it is never withheld from an engaged party.
\end{itemize}

Each hint has a primary delivery mechanism and at least one documented recovery
path: a secondary delivery mechanism (a different NPC, a geographic callback,
a later scene) that activates if the primary delivery is missed. In 3 of 7
campaigns, recovery paths were activated; in no campaign was a hint permanently
undelivered. Recovery path design is a non-negotiable requirement: the hint
structure's reliability depends on redundancy at the delivery level.

The simultaneous delivery of all four hints in the finale session is the spine's
most distinctive structural outcome. All four hints are individually available
to the party throughout the campaign; what the finale provides is the
\textit{context in which they converge}. The party does not necessarily receive
new information at the finale; they recognize what the hints have been building toward.

\subsection{Component 3: Branching End-Conditions}

Each campaign has 2--4 possible endings (Routes A through E), each with explicit
conditions. Table~\ref{tab:routes} shows the typical route structure and its
relationship to party engagement level.

\begin{table}[t]
\caption{Typical Campaign Spine Route Structure}
\label{tab:routes}
\begin{tabular}{@{}lp{4.6cm}p{2.2cm}@{}}
\toprule
Route & Condition summary & Character \\
\midrule
Route A & Party satisfies minimum structural conditions only & Institutional / political \\
Route B & Party engages 50--75\% of trust / token conditions & Personal-level \\
Route C & Default completion; no optimal path taken & Clean close \\
Route D & 90\%+ engagement \& optimal-path moment achieved & Full inversion \\
Route E & All conditions met; optimal path + unexpected completeness & Transcendent \\
\bottomrule
\end{tabular}
\end{table}

Routes A--C are always achievable without special design --- they require the
party to complete the adventure sequence and reach the finale location. Routes D
and E require sustained engagement across multiple sessions: each token earned,
each hint recognized, each spotlight moment inhabited contributes to the
route-threshold the party reaches by the finale.

The route structure prevents two common campaign failure modes: the ``nothing
happened'' ending (Route A or C ensures something always happens) and the
``impossible ideal'' ending (Routes A--C provide satisfying completion even
when D or E is not achieved). A campaign where Route A is humiliating to receive
will feel punishing to low-engagement players; a campaign where Route E requires
nothing beyond showing up will feel hollow to high-engagement players. The
route design must calibrate these thresholds deliberately.

\subsection{Component 4: PC-Spotlight Schedule}

Each PC is designated a primary spotlight adventure --- a session in which their
character arc is foregrounded, their PC-sheet grief-hook is activated, and the
party is positioned to witness their moment rather than compete with it. The
spotlight schedule distributes PC arc moments across the campaign rather than
concentrating them at the finale.

In the corpus, spotlight schedules that placed two or more PCs in the same session
correlated with lower Character Agency scores in the un-spotlighted sessions of
that session. The recommended schedule is one spotlight per sessions 1--5, with
session 6 as ensemble preparation and session 7 as the finale. This gives each
of four PCs a dedicated session in the first five, leaving the final two sessions
free for arc convergence without spotlight competition.

The spotlight schedule interacts directly with the deliverable specification:
the material accumulated in a PC's spotlight session becomes the raw content
that the deliverable is built from. A spotlight in Session 3 gives five sessions
of loading time before the finale deliverable is produced. A spotlight in Session 6
gives only one session of loading time, and the deliverable risks feeling
underprepared.

\subsection{Component 5: PC-Authored Finale Deliverable Specification}

The spine designates one or more PCs' character-sheet features as the source
for a deliverable that is accumulated across sessions and produced verbatim
at the finale. The specification includes: which PC, which sheet feature
(grief-hook, professional code, carried object, named wound), which sessions
provide the raw material, and what the deliverable's structure is (speech,
document, ritual act, silence, gesture).

This component is the most demanding to design: it requires a PC sheet with
sufficient density to constrain the deliverable's form, and a campaign with
enough material-events to load it. When both are present, the deliverable
is a natural outgrowth of the campaign rather than a scripted speech. In C1,
Grom's Silver-Ingot Confession was loaded by six sessions of artifact-destruction
rites, each adding a named shame to the ingot's material; the finale speech was
determined by what had been named, not by what was written in the module. In C2,
Thessaly's compact opening sentence was written in her own hand in the margin of
the reconstruction document during Session 6; the finale delivered what she had
already written.

Deliverables take four structural forms in the corpus:
\begin{enumerate}
  \item \textbf{Verbatim speech:} A PC speaks aloud, from their accumulated
    material, in a specific scene. Logged word-for-word.
  \item \textbf{Document delivery:} A PC produces or presents a written artifact
    (a compact, a dossier, a letter) whose content was accumulated across sessions.
  \item \textbf{Ritual act:} A PC performs an action (a forging, a keystone-setting,
    a field rite) that expresses the campaign's central question without speech.
  \item \textbf{Named gesture:} A PC performs a habitual act (cutting bread
    long-ways; holding a reliquary without opening it; leaving a commonplace
    book on a shelf) that is recognizable as completing their arc without requiring
    explication.
\end{enumerate}

All four forms were represented in the corpus. All 7 deliverables were produced
and logged.

\subsection{Operationalizing Emotional Completeness}
\label{subsec:operationalizing-completion}

The term \textit{emotional completeness} used in describing Route D and E outcomes
invites a circularity concern: these routes are classified by the session-runner's
assessment, which is informed by the same rubric used to score sessions. To reduce
circularity, we operationalize campaign completion through three rubric-external
indicators:

\begin{enumerate}
  \item \textbf{PC-authored finale deliverable produced (binary):} The deliverable
    specified in Component 5 was produced and logged verbatim in the Session 7
    log. This is a documentary check: either the verbatim text appears in the log
    or it does not.
  \item \textbf{All 4 hints delivered (binary):} Each of the four designed hints
    (primary or recovery path) was delivered before the finale. This is verifiable
    from the hint-delivery log without reference to rubric scoring.
  \item \textbf{Route D or E achieved (binary):} The party satisfied the specific
    conditions listed in the spine's end-condition document for the campaign's two
    hardest routes, which require maximal engagement with the trust/token system.
    These conditions are written in the spine before Session 1 and are observable
    independently of rubric quality scoring.
\end{enumerate}

These three indicators are observable without rubric scoring. They function as a
convergent-validity check on the Route D/E classification: a campaign that satisfies
all three has produced the spine's intended outcome by non-circular criteria.

\subsection{Component Interactions}

The five components are not independent. Their designed interactions are as
important as the components themselves:

\paragraph{Hints enable route achievement.}
Without all four hints delivered, Route D and E conditions are structurally
unachievable in most campaign designs. Hint 3 typically reveals why the surface
conflict is not the real conflict; Hint 4 reveals the specific mechanism for
resolution. Without both, the party cannot reach the optimal end-condition even
with maximal engagement.

\paragraph{Spotlights load deliverables.}
The PC-authored deliverable is built from the raw material that spotlight adventures
provide. A spotlight schedule that places spotlights in sessions 5--7 does not
provide adequate loading time for the deliverable. The optimal window is sessions 1--5.

\paragraph{End-conditions determine hint delivery priority.}
In campaigns where Route D/E requires all four hints, hint delivery reliability
at every point is non-negotiable. In campaigns where Route A or B is an acceptable
optimal, hints can have weaker primary delivery mechanisms (relying more on
recovery paths) without campaign failure.

\paragraph{Central question drives thematic coherence.}
All four hints must advance the \textit{same} question from different angles. A hint
that introduces a thematically disconnected fact --- however interesting as world-building
--- occupies a hint slot without building toward the finale convergence. The central
question functions as the thematic filter for hint design.

\begin{table}[t]
\caption{Component Interactions: Dependencies and Enables}
\label{tab:component-interaction}
\begin{tabular}{@{}p{3.2cm}p{4.6cm}@{}}
\toprule
Component & Enables \\
\midrule
Hints 3 and 4 & Route D and E conditions \\
Spotlight schedule (S1--S5) & PC-authored deliverable loading \\
End-condition routes & Hint delivery priority and recovery path design \\
Central question & Hint thematic coherence \\
PC-authored deliverable spec & Session-to-session continuity toward a single endpoint \\
\bottomrule
\end{tabular}
\end{table}
